\section{Using and reproducing the implementations}
\label{sec:sw-reproduce}

\subsection{A fixed identity metric for a TensorFlow model}

The following small example illustrates the explicit tensor interface. Its
short fixed counts are a mechanics example, not a posterior protocol. The
Gaussian log density can be replaced by an independent-row TensorFlow model
or by the frozen transformed density in equation~\eqref{eq:sw-neutra-target}.
The factory supplies stable signatures and XLA compilation; the caller must
choose a device, keep data placement explicit, and check validity. GPU use
requires a compatible installed TensorFlow/CUDA environment.

\begin{lstlisting}[language=Python,basicstyle=\ttfamily\scriptsize,breaklines=true]
import os
os.environ['TF_FORCE_GPU_ALLOW_GROWTH'] = 'true'
import tensorflow as tf
gpus = tf.config.list_physical_devices('GPU')
if not gpus:
    raise RuntimeError('This example requires a GPU')
for gpu in gpus:
    tf.config.experimental.set_memory_growth(gpu, True)
    assert tf.config.experimental.get_memory_growth(gpu)
from docs.benchmarks.walnutpie_gpu_tf import (
    Config, make_supplied_mass_functions)

def log_density(q, location):
    return -.5 * tf.reduce_sum((q-location)**2, axis=-1)

cfg = Config(chains=4, dimension=10, warmup=20, draws=20,
             adapt_mass=False)
f = make_supplied_mass_functions(tf, log_density, cfg,
    model_signature=[tf.TensorSpec([10], tf.float64)])
with tf.device('/GPU:0'):
    q = tf.ones([4, 10], tf.float64) * .2
    keys = tf.constant([[101, 0], [102, 0],
                        [103, 0], [104, 0]], tf.int32)
    mass = tf.ones_like(q)
    epsilon = tf.fill([4], tf.constant(.1, tf.float64))
    fit = f.fit(q, keys, mass, epsilon,
                tf.zeros([10], tf.float64))
assert bool(tf.reduce_all(fit.valid).numpy())
tf.debugging.assert_equal(fit.warmup.state.inverse_mass, mass)
retained = fit.sampling.draws  # [20, 4, 10]; warmup excluded
\end{lstlisting}

Here \texttt{adapt\_mass=False} freezes the supplied momentum covariance,
but step-size and minimum-work adaptation still operate during warmup.
The factory's \texttt{adapt\_step=False} option skips only the initial
step search; it does not disable warmup adaptation. To use already frozen
R2P tuning, set zero warmup, provide the step explicitly, skip that search,
and invoke sampling on the initialized state. A change of model data or
target requires reinitialization of density/score caches and reassessment
of tuning. Continuing an unchanged sampling state preserves the RNG and
frozen tuning.

The numerical state contains position, density, score, random keys and
counters, step size, inverse mass, minimum work, adaptation iteration,
Adam accumulators and discounted moments. Dense initialization additionally
stores the fixed basis. Run-control checkpoints also contain iteration
counts and log-density moments. Save this complete state for continuation;
saving only the last position does not reproduce the run. Fixed-capacity
chunk outputs carry a mask, and padded rows must never enter estimates.

\subsection{Tuning and comparison protocol}

For the standalone Stan-style route, first validate the unconstrained
target and score, Jacobians, row isolation, finite starts and metric
convention. Use the inspected Stan three-stage warmup, its metric-boundary
step search and dual-averaging restart, then freeze adaptation before
retaining samples. The adaptation-only wrapper rejects unsupported short
budgets; the direct port separately preserves the pinned Stan short-budget
fallback described in Section~\ref{sec:sw-stan-tuning}. These are different
interface contracts. Source defaults nominate a starting configuration;
they do not certify that a target is adequately explored.

For Walnutpie, record the deterministic or R2P variant, base interval,
refinement bounds, jitter, acceptance statistic, mass convention and
initialization search limit. Its default Adam warmup is not Stan dual
averaging, and R2P combined with that warmup is not the paper's Appendix C
quantile algorithm. The author-replication driver runs Appendix C in the
independent Python reference and transfers its frozen parameters into XLA.
An upstream log-density stopping statistic or step/mass agreement can
finish adaptation without establishing parameter convergence.

Compare tuned fixed-length HMC with a declared search over step size and
leapfrog count, charging its calibration and selection work. Compare
NUTS with its complete adaptation and charge that work too. Use held-out
retained streams for the posterior comparison. Report cold compilation,
warmup, retained execution, diagnostics and total elapsed time separately;
give the exact number of chains, device, dtype, model/gradient backend and
work-counter meaning. For execution-layout comparisons, use common inputs,
balanced order, synchronization and uncertainty over independent seed blocks.
For different samplers, assess validity and posterior agreement before
ranking effective information per unit time. Shared seeds do not justify
an independent-stream standard-error formula.

\subsection{Source and test map}

The authoritative implementation and test locations, relative to the
repository root, are listed below. Reported historical test totals refer
to their recorded revisions, not a promise that a future suite has the
same test count. Native-header tests require the pinned local sources;
missing dependencies must not silently substitute a different reference.

{\small
\begin{longtable}{>{\raggedright\arraybackslash}p{.24\linewidth}>{\raggedright\arraybackslash}p{.68\linewidth}}
\toprule
Component & Implementation and evidence\\
\midrule
\endhead
Stan warmup & \path{docs/benchmarks/stan_nuts_warmup.py};
\path{docs/reference/nuts-tuning.md}; adaptation code-audit result.\\
Direct Stan and TFP & \path{docs/benchmarks/stan_nuts.py},
\path{docs/benchmarks/stan_nuts_tfp.py};
\path{tests/test_stan_nuts_tfp.py}; direct-port and native-TFP result notes.\\
Native bridge & \path{docs/benchmarks/native_sampler_bridge/};
native Stan/Walnutpie result note and source manifests.\\
Stan GPU models & \path{docs/benchmarks/stan_gpu_tf/};
\path{tests/test_stan_gpu_tf.py}.\\
Walnutpie numerical core & \path{docs/benchmarks/walnutpie_gpu_tf/core.py},
\path{api.py}, \path{r2p.py} in that package;
\path{tests/test_walnutpie_gpu_tf.py}.\\
Initialization, control, mass & \path{initialization.py},
\path{control.py}, \path{mass_coordinates.py} in the same package;
\path{tests/test_walnutpie_gap_closure.py},
\path{tests/test_walnutpie_supplied_mass.py}.\\
Funnel and hybrid studies & \path{author_reference/},
\path{funnel_master.py}, \path{replication.py},
\path{replication_analysis.py}, \path{hybrid_benchmark.py} in that package;
\path{tests/test_walnutpie_hybrid.py}.\\
Full80 application & \path{docs/benchmarks/full80_walnutpie_sampling.py},
its support, controller-helper and report modules;
\path{tests/test_full80_walnutpie.py}.\\
\bottomrule
\end{longtable}}

The Walnutpie C++ oracle is also saved under
\path{docs/benchmarks/walnutpie_gpu_tf/prebuilt/linux-x86_64/20261029/}.
Its executable is 249,008 bytes (243.2 KiB); the matching source archive
is about 1 MiB. Checksums, compiler flags, pinned headers, Eigen 3.4.0 and
licences accompany it. It requires Linux x86-64, glibc 2.34 or newer and
libstdc++ symbols through \texttt{GLIBCXX\_3.4.29}/\texttt{CXXABI\_1.3.9}.
It has no TensorFlow or CUDA dependency. It emits 528 deterministic JSON
oracle records and implements the supplied analytic targets, not arbitrary
user TensorFlow models. The recorded generic build took 7.33 seconds.
Platform-specific compilation remains necessary elsewhere, and each new
TensorFlow model signature still requires XLA compilation.

For a bounded CPU correctness check with the prebuilt reference, use a fresh
output directory and the installed TensorFlow environment:
\begin{lstlisting}[basicstyle=\ttfamily\scriptsize,breaklines=true]
CUDA_VISIBLE_DEVICES=-1 TF_FORCE_GPU_ALLOW_GROWTH=true \
python -m docs.benchmarks.walnutpie_gpu_tf.worker \
  --mode validate --device cpu --execution-layout packed \
  --oracle docs/benchmarks/walnutpie_gpu_tf/prebuilt/linux-x86_64/20261029/oracle \
  --output /tmp/walnutpie-validation-new

CUDA_VISIBLE_DEVICES=-1 TF_FORCE_GPU_ALLOW_GROWTH=true \
python -m pytest -q tests/test_full80_walnutpie.py
\end{lstlisting}

Do not pass the tracked executable as a build destination: the original
master's build, qualify and all phases rebuild that path. Worker validation
or the master's validate/benchmark phase can consume the prebuilt directly.
Full GPU benchmarks and million-iteration replications are separate,
budgeted experiments; their launch commands, source pins and stop conditions
belong to the saved plans, rather than this small correctness example.

\subsection{Result records and reproducibility limits}

The result-note index below gives filenames under \path{docs/plans/}.
Each note links its run manifests, exact commands, seeds, environment,
timing boundaries and structured diagnostics. The monograph evidence
snapshot under \path{docs/reference/stan-walnutpie-evidence-20261031/}
preserves selected completed-result JSON files, their original locations
and SHA-256 hashes. Large raw draws, machine-local source installations
and financial-model inputs remain at the recorded experiment locations;
the compact snapshot is not a replacement for those inputs.

{\small
\begin{longtable}{>{\raggedright\arraybackslash}p{.23\linewidth}>{\raggedright\arraybackslash}p{.69\linewidth}}
\toprule
Study & Result note or completed record\\
\midrule
\endhead
Early devices/HMC & \path{bayesfilter-nuts-xla-benchmark-result-2026-10-27.md};
\path{bayesfilter-nuts-fixed-hmc-result-2026-10-27.md};
\path{bayesfilter-nuts-geometry-overhead-result-2026-10-27.md}.\\
Stan transfer/audit & \path{bayesfilter-stan-nuts-transplant-result-2026-10-27.md};
\path{bayesfilter-stan-adaptation-code-audit-result-2026-10-27.md};
\path{bayesfilter-full-stan-nuts-port-result-2026-10-27.md}.\\
TFP and GPU Stan & \path{bayesfilter-stan-tfp-native-result-2026-10-28.md};
\path{bayesfilter-native-stan-walnutpie-result-2026-10-28.md};
\path{bayesfilter-stan-gpu-tf-result-2026-10-28.md}.\\
Walnutpie transfer & \path{bayesfilter-walnutpie-gpu-tf-result-2026-10-28.md};
\path{bayesfilter-walnutpie-gap-closure-result-2026-10-29.md};
\path{bayesfilter-walnutpie-supplied-mass-result-2026-10-29.md}.\\
Optimization/hybrid & \path{bayesfilter-walnutpie-optimization-result-2026-10-29.md};
\path{bayesfilter-walnutpie-hybrid-large-model-result-2026-10-29.md}.\\
Author funnel & \path{bayesfilter-walnuts-author-funnel-result-2026-10-29.md};
million-run \path{REPORT.md} and \path{comparison.json} in
\path{docs/benchmarks/artifacts/walnuts-million-analytic-20261029-01/}.\\
q20 & \path{bayesfilter-q20-walnutpie-placement-results-2026-10-29.md};
\path{bayesfilter-q20-walnutpie-estimation-results-2026-10-29.md},
including its 31 October terminal update.\\
Full80 & \path{bayesfilter-full80-walnutpie-result-2026-10-30.md};
\path{comparison-01/} and \path{terminal-review-02/} under
\path{docs/benchmarks/artifacts/full80-walnutpie-20261030/}.\\
\bottomrule
\end{longtable}}

The remaining limits are specific. Appendix C tuning has not been
transferred to a native TensorFlow kernel. The exact extreme-start
$v=-30$ transient study is incomplete. Dense supplied-mass adaptation is
diagonal in a fixed basis. The C++ bridge does not fuse arbitrary TF
models, and the pure TensorFlow port does not reproduce upstream random
streams bit for bit. CPU/GPU preference depends on model cost, batching
and compiler behaviour. The full80 saved-draw comparison is complete, but
fresh target replay requires recovery of its bound MacroFinance source
closure; q20 still lacks an accepted posterior estimate. These distinctions
identify what further evidence is needed without discarding the completed
source checks, execution measurements or posterior comparisons.
